\documentclass{amsart}
% For dealing with references we use the comment environment
\usepackage{verbatim}
\newenvironment{reference}{\comment}{\endcomment}
%\newenvironment{reference}{}{}
\newenvironment{slogan}{\comment}{\endcomment}
\newenvironment{history}{\comment}{\endcomment}

% For commutative diagrams we use Xy-pic
\usepackage[all]{xy}

% We use 2cell for 2-commutative diagrams.
\xyoption{2cell}
\UseAllTwocells

% We use multicol for the list of chapters between chapters
\usepackage{multicol}

% This is generally recommended for better output
\usepackage{lmodern}
\usepackage[T1]{fontenc}

% For cross-file-references

% Package for hypertext links:
\usepackage{hyperref}

% For any local file, say "hello.tex" you want to link to please

% Theorem environments.
%
\theoremstyle{plain}
\newtheorem{theorem}[subsection]{Theorem}
\newtheorem{proposition}[subsection]{Proposition}
\newtheorem{lemma}[subsection]{Lemma}

\theoremstyle{definition}
\newtheorem{definition}[subsection]{Definition}
\newtheorem{example}[subsection]{Example}
\newtheorem{exercise}[subsection]{Exercise}
\newtheorem{situation}[subsection]{Situation}

\theoremstyle{remark}
\newtheorem{remark}[subsection]{Remark}
\newtheorem{remarks}[subsection]{Remarks}

\numberwithin{equation}{subsection}

% Macros
%
\def\lim{\mathop{\mathrm{lim}}\nolimits}
\def\colim{\mathop{\mathrm{colim}}\nolimits}
\def\Spec{\mathop{\mathrm{Spec}}}
\def\Hom{\mathop{\mathrm{Hom}}\nolimits}
\def\Ext{\mathop{\mathrm{Ext}}\nolimits}
\def\SheafHom{\mathop{\mathcal{H}\!\mathit{om}}\nolimits}
\def\SheafExt{\mathop{\mathcal{E}\!\mathit{xt}}\nolimits}
\def\Sch{\mathit{Sch}}
\def\Mor{\mathop{\mathrm{Mor}}\nolimits}
\def\Ob{\mathop{\mathrm{Ob}}\nolimits}
\def\Sh{\mathop{\mathit{Sh}}\nolimits}
\def\NL{\mathop{N\!L}\nolimits}
\def\CH{\mathop{\mathrm{CH}}\nolimits}
\def\proetale{{pro\text{-}\acute{e}tale}}
\def\etale{{\acute{e}tale}}
\def\QCoh{\mathit{QCoh}}
\def\Ker{\mathop{\mathrm{Ker}}}
\def\Im{\mathop{\mathrm{Im}}}
\def\Coker{\mathop{\mathrm{Coker}}}
\def\Coim{\mathop{\mathrm{Coim}}}

% Boxtimes
%
\DeclareMathSymbol{\boxtimes}{\mathbin}{AMSa}{"02}

%
% Macros for moduli stacks/spaces
%
\def\QCohstack{\mathcal{QC}\!\mathit{oh}}
\def\Cohstack{\mathcal{C}\!\mathit{oh}}
\def\Spacesstack{\mathcal{S}\!\mathit{paces}}
\def\Quotfunctor{\mathrm{Quot}}
\def\Hilbfunctor{\mathrm{Hilb}}
\def\Curvesstack{\mathcal{C}\!\mathit{urves}}
\def\Polarizedstack{\mathcal{P}\!\mathit{olarized}}
\def\Complexesstack{\mathcal{C}\!\mathit{omplexes}}
% \Pic is the operator that assigns to X its picard group, usage \Pic(X)
% \Picardstack_{X/B} denotes the Picard stack of X over B
% \Picardfunctor_{X/B} denotes the Picard functor of X over B
\def\Pic{\mathop{\mathrm{Pic}}\nolimits}
\def\Picardstack{\mathcal{P}\!\mathit{ic}}
\def\Picardfunctor{\mathrm{Pic}}
\def\Deformationcategory{\mathcal{D}\!\mathit{ef}}

\setlength{\emergencystretch}{3em}
\begin{document}
\title[Stacks Project, Tag 0GAA]{358 --- Uniform twist vanishing for almost-projective modules on a blowup}
\maketitle
\noindent Copyright (C) 2005--2025 Johan de Jong. Stacks Project authors. Source: \href{https://stacks.math.columbia.edu/tag/0GAA}{Tag 0GAA}.

\noindent Editorial difficulty note (not author text). Difficulty H2: The proof splits the failure of pullback exactness into two coherent short exact sequences with error supported on a controlled exceptional thickening. Descending induction on cohomological degree builds explicit twist shifts that annihilate transition maps, and a separate boundary-map argument controls degree-zero images. The constants depend on the associated graded ring and error thickness, so ordinary Serre vanishing does not establish the conclusion..

\noindent Editorial provenance note (not author text). History: excerpt from revision \texttt{a0444\allowbreak 6e57e\allowbreak c1fbc\allowbreak 252a8\allowbreak 71afc\allowbreak ec775\allowbreak 2fb28\allowbreak 07b14}, local-cohomology.tex, lines 4927--5073. Reference presentation was adapted; original-author context and core support blocks were retained or added. The mathematical wording is unchanged. Licensed under GNU FDL 1.2 or later; no invariant sections or cover texts. See ../licenses/COPYING. Context blocks reproduce author definitions and supporting results; they are not additional counted problems.

\section{A bit of uniformity, III}
\label{section-uniform}

\noindent
In this section we fix a Noetherian ring $A$ and an ideal $I \subset A$.
Our goal is to prove Lemma \href{https://stacks.math.columbia.edu/tag/0GAD}{lemma bound two term complex (Tag 0GAD)} which we will
use in a later chapter to solve a lifting problem, see
Algebraization of Formal Spaces, Lemma
\href{https://stacks.math.columbia.edu/tag/0GAQ}{lemma get morphism general better (Tag 0GAQ)}.

\medskip\noindent
Throughout this section we denote
$$
p : X \to \Spec(A)
$$
the blowing up of $\Spec(A)$ in the ideal $I$. In other words, $X$ is the
$\text{Proj}$ of the Rees algebra $\bigoplus_{n \geq 0} I^n$. We also consider
the fibre product
$$
\xymatrix{
Y \ar[r] \ar[d] & X \ar[d]^p \\
\Spec(A/I) \ar[r] & \Spec(A)
}
$$
Then $Y$ is the exceptional divisor of the blowup and hence
an effective Cartier divisor on $X$ such that
$\mathcal{O}_X(-1) = \mathcal{O}_X(Y)$. Since taking $\text{Proj}$
commutes with base change we have
$$
Y = \text{Proj}(\bigoplus\nolimits_{n \geq 0} I^n/I^{n + 1}) = \text{Proj}(S)
$$
where $S = \text{Gr}_I(A) = \bigoplus_{n \geq 0} I^n/I^{n + 1}$.

\medskip\noindent
We denote
$d = d(S) = d(\text{Gr}_I(A)) = d(\bigoplus_{n \geq 0} I^n/I^{n + 1})$
the maximum of the dimensions of the fibres of $p$
(and we set it equal to $0$ if $X = \emptyset$).
This is well defined. In fact, we have
\begin{enumerate}
\item $d \leq t - 1$ if $I = (a_1, \ldots, a_t)$ since then
$X \subset \mathbf{P}^{t - 1}_A$, and
\item $d$ is also the maximal dimension of the fibres of
$\text{Proj}(S) \to \Spec(S_0)$ provided that $Y$
is nonempty and $d = 0$ if $Y = \emptyset$
(equivalently $S = 0$, equivalently $I = A$).
\end{enumerate}
Hence $d$ only depends on the isomorphism class of $S = \text{Gr}_I(A)$.
Observe that $H^i(X, \mathcal{F}) = 0$ for every coherent
$\mathcal{O}_X$-module $\mathcal{F}$ and $i > d$ by
Cohomology of Schemes, Lemmas
\href{https://stacks.math.columbia.edu/tag/02V7}{lemma higher direct images zero above dimension fibre (Tag 02V7)} and
\href{https://stacks.math.columbia.edu/tag/01XK}{lemma quasi coherence higher direct images application (Tag 01XK)}.
Of course the same is true for coherent modules on $Y$.

\medskip\noindent
We denote
$q = q(S) = q(\text{Gr}_I(A)) = q(\bigoplus_{n \geq 0} I^n/I^{n + 1})$
the integer defined as follows. Note that the algebra
$S = \bigoplus_{n \geq 0} I^n/I^{n + 1}$
is a Noetherian graded ring generated in degree $1$ over degree $0$.
Hence by
Cohomology of Schemes, Lemmas \href{https://stacks.math.columbia.edu/tag/0AG6}{lemma coherent on proj (Tag 0AG6)} and
\href{https://stacks.math.columbia.edu/tag/0AG7}{lemma recover tail graded module (Tag 0AG7)} we can define $q(S)$
as the smallest integer $q(S) \geq 0$ such that for all $q \geq q(S)$ we have
$H^i(Y, \mathcal{O}_Y(q)) = 0$ for $1 \leq i \leq d$ and
$H^0(Y, \mathcal{O}_Y(q)) = I^q/I^{q + 1}$.
(If $S = 0$, then $q(S) = 0$.)

\medskip\noindent
For $n \geq 1$ we may consider the effective Cartier divisor $nY$
which we will denote $Y_n$.



\section{Modules which are close to being projective}
\label{more-algebra-section-near-projective}

\noindent
There seem to be many different of definitions in the literature of
``almost projective modules''. In this section we discuss just one
of the many possibilities.



\begin{definition}
\label{more-algebra-definition-near-projective}
Let $R$ be a ring. Let $I \subset R$ be an ideal. Let $M$ be an $R$-module.
We say $M$ is {\it $I$-projective}\footnote{This is nonstandard notation.}
if the equivalent conditions of Lemma \ref{more-algebra-lemma-near-projective} hold.
\end{definition}

\begin{lemma}
\label{more-algebra-lemma-near-projective}
Let $R$ be a ring. Let $I \subset R$ be an ideal. Let $M$ be an $R$-module.
The following conditions are equivalent
\begin{enumerate}
\item for every $a \in I$ the map $a : M \to M$ factors through a projective
$R$-module,
\item for every $a \in I$ the map $a : M \to M$ factors through a free
$R$-module, and
\item $\Ext^1_R(M, N)$ is annihilated by $I$ for every $R$-module $N$.
\end{enumerate}
\end{lemma}

\noindent
Let us introduce a notation: given $n \geq c \geq 0$
{\it an $(A, n, c)$-module}
is a finite $A$-module $M$ which is annihilated by $I^n$
and which as an $A/I^n$-module is $I^c/I^n$-projective, see
More on Algebra, Section \ref{more-algebra-section-near-projective}.

\medskip\noindent
We will use the following abuse of notation: given an $A$-module $M$
we denote $p^*M$ the quasi-coherent module gotten by pulling back by $p$
the quasi-coherent module $\widetilde{M}$ on $\Spec(A)$ associated to $M$.
For example we have $\mathcal{O}_{Y_n} = p^*(A/I^n)$.
For a short exact sequence $0 \to K \to L \to M \to 0$ of $A$-modules
we obtain an exact sequence
$$
p^*K \to p^*L \to p^*M \to 0
$$
as $\widetilde{\ }$ is an exact functor and $p^*$ is a right exact functor.


\begin{lemma}
\label{lemma-vanishing-coh-almost-projective}
With $q_0 = q(S)$ and $d = d(S)$ as above, suppose we have
integers $n \geq c \geq 0$, an $(A, n, c)$-module $M$,
an index $i \in \{0, 1, \ldots, d\}$, and an integer $q$.
Then we distinguish the following cases
\begin{enumerate}
\item In the case $i = d \geq 1$ and $q \geq q_0$ we have
$H^d(X, p^*M(q)) = 0$.
\item In the case $i = d - 1 \geq 1$ and $q \geq q_0$ we have
$H^{d - 1}(X, p^*M(q)) = 0$.
\item In the case $d - 1 > i > 0$ and $q \geq q_0 + (d - 1 - i)c$
the map
$H^i(X, p^*M(q)) \to H^i(X, p^*M(q - (d - 1 - i)c))$
is zero.
\item In the case $i = 0$, $d \in \{0, 1\}$, and $q \geq q_0$, there
is a surjection
$$
I^qM \longrightarrow H^0(X, p^*M(q))
$$
\item In the case $i = 0$, $d > 1$, and $q \geq q_0 + (d - 1)c$ the map
$$
H^0(X, p^*M(q)) \to H^0(X, p^*M(q - (d - 1)c))
$$
has image contained in the image of the canonical map
$I^{q - (d - 1)c}M \to H^0(X, p^*M(q - (d - 1)c))$.
\end{enumerate}
\end{lemma}

\begin{proof}
Let $M$ be an $(A, n, c)$-module. Choose a short exact sequence
$$
0 \to K \to (A/I^n)^{\oplus r} \to M \to 0
$$
We will use below that $K$ is an $(A, n, c)$-module, see More on Algebra,
Lemma \href{https://stacks.math.columbia.edu/tag/0G95}{lemma ses near projective (Tag 0G95)}.
Consider the corresponding exact sequence
$$
p^*K \to (\mathcal{O}_{Y_n})^{\oplus r} \to p^*M \to 0
$$
We split this into short exact sequences
$$
0 \to \mathcal{F} \to p^*K \to \mathcal{G} \to 0
\quad\text{and}\quad
0 \to \mathcal{G} \to (\mathcal{O}_{Y_n})^{\oplus r} \to p^*M \to 0
$$
By Lemma \href{https://stacks.math.columbia.edu/tag/0GA8}{lemma almost exactness (Tag 0GA8)} the coherent module $\mathcal{F}$
is scheme theoretically supported on $Y_c$.

\medskip\noindent
Proof of (1). Assume $d > 0$. We have to prove
$H^d(X, p^*M(q)) = 0$ for $q \geq q_0$.
By the vanishing of the cohomology of twists of $\mathcal{G}$ in degrees $> d$
and the long exact cohomology sequence associated to the second
short exact sequence above, it suffices to prove that
$H^d(X, \mathcal{O}_{Y_n}(q)) = 0$.
This is true by Lemma \href{https://stacks.math.columbia.edu/tag/0GA7}{lemma bound q and d (Tag 0GA7)}.

\medskip\noindent
Proof of (2). Assume $d > 1$. We have to prove
$H^{d - 1}(X, p^*M(q)) = 0$ for $q \geq q_0$.
Arguing as in the previous paragraph, we see that it suffices
to show that $H^d(X, \mathcal{G}(q)) = 0$. Using the first
short exact sequence and the vanishing of the cohomology
of twists of $\mathcal{F}$ in degrees $> d$ we see that it suffices
to show $H^d(X, p^*K(q))$ is zero which is
true by (1) and the fact that $K$ is an $(A, n, c)$-module (see above).

\medskip\noindent
Proof of (3). Let $0 < i < d - 1$ and assume the statement holds for $i + 1$
except in the case $i = d - 2$ we have statement (2).
Using the long exact sequence of cohomology associated to the second
short exact sequence above we find an injection
$$
H^i(X, p^*M(q - (d - 1 - i)c)) \subset
H^{i + 1}(X, \mathcal{G}(q - (d - 1 - i)c))
$$
as $q - (d - 1 - i)c \geq q_0$ gives the vanishing of
$H^i(X, \mathcal{O}_{Y_n}(q - (d - 1 - i)c))$
(see above). Thus it suffices to show that the map
$H^{i + 1}(X, \mathcal{G}(q)) \to H^{i + 1}(X, \mathcal{G}(q - (d - 1 - i)c))$
is zero. To study this, we consider the maps of exact sequences
$$
\xymatrix{
H^{i + 1}(X, p^*K(q)) \ar[r] \ar[d] &
H^{i + 1}(X, \mathcal{G}(q)) \ar[r] \ar[d] \ar@{..>}[ld] &
H^{i + 2}(X, \mathcal{F}(q)) \ar[d] \\
H^{i + 1}(X, p^*K(q - c)) \ar[r] \ar[d] &
H^{i + 1}(X, \mathcal{G}(q - c)) \ar[r] \ar[d] &
H^{i + 2}(X, \mathcal{F}(q - c)) \\
H^{i + 1}(X, p^*K(q - (d - 1 - i)c)) \ar[r] &
H^{i + 1}(X, \mathcal{G}(q - (d - 1 - i)c))
}
$$
Since $\mathcal{F}$ is scheme theoretically supported on $Y_c$
we see that the canonical map
$\mathcal{G}(q) \to \mathcal{G}(q - c)$ factors through
$p^*K(q - c)$ by Lemma \href{https://stacks.math.columbia.edu/tag/0GA9}{lemma annihilated (Tag 0GA9)}.
This gives the dotted arrow in the diagram. (In fact, for the proof it
suffices to observe that the vertical arrow on the extreme right is
zero in order to get the dotted arrow as a map of sets.)
Thus it suffices to show that
$H^{i + 1}(X, p^*K(q - c)) \to
H^{i + 1}(X, p^*K(q - (d - 1 - i)c))$
is zero. If $i = d - 2$, then the source of this arrow
is zero by (2) as $q - c \geq q_0$ and $K$ is an $(A, n, c)$-module.
If $i < d - 2$, then as $K$ is an $(A, n, c)$-module, we get from the
induction hypothesis that the map is indeed zero
since $q - c - (q - (d - 1 - i)c) = (d - 2 - i)c = (d - 1 - (i + 1))c$
and since $q - c \geq q_0 + (d - 1 - (i + 1))c$.
In this way we conclude the proof of (3).

\medskip\noindent
Proof of (4). Assume $d \in \{0, 1\}$ and $q \geq q_0$.
Then the first short exact sequence gives a surjection
$H^1(X, p^*K(q)) \to H^1(X, \mathcal{G}(q))$
and the source of this arrow is zero by case (1). Hence
for all $q \in \mathbf{Z}$ we see that the map
$$
H^0(X, (\mathcal{O}_{Y_n})^{\oplus r}(q))
\longrightarrow
H^0(X, p^*M(q))
$$
is surjective. For $q \geq q_0$ the
source is equal to $(I^q/I^{q + n})^{\oplus r}$ by
Lemma \href{https://stacks.math.columbia.edu/tag/0GA7}{lemma bound q and d (Tag 0GA7)} and this easily proves the statement.

\medskip\noindent
Proof of (5). Assume $d > 1$. Arguing as in the proof of (4) we see that
it suffices to show that the image of
$$
H^0(X, p^*M(q))
\longrightarrow
H^0(X, p^*M(q - (d - 1)c))
$$
is contained in the image of
$$
H^0(X, (\mathcal{O}_{Y_n})^{\oplus r}(q - (d - 1)c))
\longrightarrow
H^0(X, p^*M(q - (d - 1)c))
$$
To show the inclusion above, it suffices to show that for
$\sigma \in H^0(X, p^*M(q))$ with boundary
$\xi \in H^1(X, \mathcal{G}(q))$ the image of $\xi$ in
$H^1(X, \mathcal{G}(q - (d - 1)c))$ is zero. This follows by the
exact same arguments as in the proof of (3).
\end{proof}
\end{document}
